\section{Introduction}
Expander graphs are known to represent a class of easier instances for many problems. Therefore, breaking down the input into disjoint expanders can allow to conveniently solve the task on each of them separately, before combining the partial results into a global solution. This approach is enabled by $(\epsilon, \phi)$-expander decompositions (for short, $(\epsilon, \phi)$-ED). For an undirected graph $G=(V,E)$, this is a partition $\mathcal{U}$ of the vertex set $V$, such that there are at most $\epsilon |E|$ inter-cluster edges, while every cluster $U \in \mathcal{U}$ induces a $\phi$-expander. The list of successful applications of this framework is long, including Laplacian system solvers~\cite{stlapl}, deterministic algorithms for minimum cut~\cite{mincutjli}, graph and hyper-graph sparsification~\cite{stspectral, chugao, hgsparse,svapprox}, dynamic algorithms for cut and connectivity problems~\cite{boundlink,sublinearconn}, fast max flow algorithms~\cite{maxflow}, distributed triangle enumeration~\cite{triang}, polynomial time algorithms for semirandom planted CSPs~\cite{csps}, and many more.

We refer to $\epsilon$ and $\phi$ as the \emph{sparsity} parameters: the former controls how sparsely connected the clusters need to be, and the latter determines how expanding (i.e. non-sparse) the cuts within clusters are. The reason for using the same term for both is that they are in fact very closely related. One can show that any $n$-vertex graph has an $(\epsilon,\phi)$-ED with $\epsilon = O(\phi \log n)$. To see this, consider the following constructive argument: if the graph has no $\phi$-sparse cut then $\{V\}$ is a valid ED of~$G$, otherwise recurse on the two sides $(S,V \setminus S)$ of a $\phi$-sparse cut and union the results to get an ED for $G$. Every cluster in this decomposition is then an expander, and a charging argument allows to bound the number of inter-cluster edges by $O(\phi |E| \log n)$~\cite{offlineexpdec}. One can also observe that no better asymptotic trade-off between $\epsilon$ and $\phi$ is possible in general~\cite{clusteringhypercube}. Therefore, this sets the benchmark for algorithmic constructions of EDs.

At a high level, many ED algorithms follow the approach suggested by the existential argument. As a naive implementation would take exponential time, the crux often lies in efficient algorithms that either certify that a large portion of the input is an expander, or find a balanced sparse cut. This would result in a small depth recursion, thanks to balancedness, where each level requires little computational resources. There are several successful examples of this approach. In the sequential setting, a recent algorithm constructs a $(O(\phi \log^3 n), \phi)$-ED in $\otil(|E|)$ time~\cite{liexpdec}, based on the previous best algorithm which runs in $\otil(|E|/\phi)$ time~\cite{offlineexpdec}. There is also a deterministic counterpart, which outputs a $(\phi \cdot n^{o(1)}, \phi)$-ED in almost linear time~\cite{determexpdec}. For the CONGEST model of distributed computing, it is possible to obtain, for instance, a $(\phi^{1/\sqrt{\log n}} \cdot n^{o(1)}, \phi)$-ED in $n^{o(1)} / \phi$ rounds~\cite{triang}. It is also possible to maintain a $(\phi \cdot n^{o(1)},\phi)$-ED for a graph undergoing edge updates in $n^{o(1)}/\phi^2$ amortized update time~\cite{boundlink,dynexpdec}.

\subsection{Previous work}
In the streaming setting, the problem of finding EDs was open until the recent work of~\cite{streamexpdec}. They obtain a dynamic stream algorithm which outputs a $(O(\phi \log n),\phi)$-ED and takes $\otil(n/\phi^2)$ space. While being optimal in the quality of the decomposition, decoding the sketch to actually output an ED takes exponential time. The authors also give a polynomial time version: one can produce a $(\phi \cdot n^{o(1)}, \phi)$-ED using space $\otil(n/\phi^2)+n^{1+o(1)}/\phi^{1-o(1)}$, with a post-processing that takes $\poly(n)$ time (where the $o(1)$'s can be tuned, allowing for a quality-space trade-off).

These streaming algorithms also adopt the recursive approach based on finding balanced $\phi$-sparse cuts, but the streaming model poses a challenge that we now illustrate. Sparsification for graph streams has been extensively studied~\cite{ahnguha,linmeas,sparsdynamicstreams,focsspars,spectralspars}, so a natural attempt would consist of maintaining a cut sparsifier as the stream comes and later run the recursive partitioning on it. However, it must be noted that these cuts are to be found in the subgraphs induced by the two sides of a previously made cut. This is not a problem in a classical computational setting, but it actually constitutes the main obstacle for the streaming model: at sketching time, the algorithm does not know which subgraphs it will need to access. Unfortunately, it is impossible to preserve cut sizes in arbitrary subgraphs with multiplicative precision. The natural work-around is to introduce an additive error term: in~\cite{streamexpdec}, the authors introduce the concept of {power-cut} sparsifiers. For any given partition $\mathcal{U}$ of $V$, these sparsifiers preserve the cuts $|E(S,U\setminus S)|$ of each cluster $U \in \mathcal{U}$ in the partition up to an error of $\delta \cdot |E(S,U\setminus S)| + \psi \cdot \vol(S)$. They further show that maintaining $\otil(n / \delta \psi)$ random linear measurements of the incidence matrix of the input graph is enough to obtain such sparsifiers. One can see that setting $\delta \ll 1$ and~$\psi \approx \phi$ preserves the sparsity of cuts to within an additive error of roughly $\phi$, using $\otil(n/\phi)$ linear measurements. This is enough to find a balanced $\phi$-sparse cut in every subgraph induced by a given partition of the vertex set. There is another caveat, though: power-cut sparsifiers give a high probability guarantee for any fixed partition, but we cannot expect them to work for all partitions simultaneously. Therefore, in order not to use the sparsifiers adaptively, we need one power-cut sparsifier for every recursion level. The authors show that the depth of the procedure cannot exceed $\otil(1/\epsilon) = \otil(1/\phi)$, thus obtaining the space complexity stated before. A few more details are involved in the polynomial time algorithm, but the underlying framework is the same.

\subsection{Our contribution}
The work of~\cite{streamexpdec} initiated the study of expander decompositions in the streaming setting, and consequently raised the question of whether a dependence in space on the sparsity parameter~$\phi$ is inherent. In this paper, we move towards settling the streaming complexity of expander decompositions by attacking the problem on two fronts: (1) we give a nearly optimal algorithm for ``one-level'' expander decomposition that avoids the sparsity dependence, and (2) we show that computing a ``repeated'' expander decomposition, commonly used in applications, cannot avoid such dependence.
\subsubsection*{Upper bound}
We give an $\otil(n)$ space algorithm for computing a $(O(\phi \log n), \phi)$-ED in dynamic streams. Specifically, we show that a ``universal'' sketch consisting of $\otil(n)$ random linear measurements of the incidence matrix can be decoded into a $(O(\phi \log n), \phi)$-ED for any $\phi$: the sketch is independent of the sparsity~$\phi$.

\begin{mainthsimple}[ED algorithm -- exponential time decoding]
	\label{cor:main}
	Let $G=(V,E)$ be a graph given in a dynamic stream. Then, there is an algorithm that maintains a linear sketch of $G$ in $\otil(n)$ space. For any $\phi \in (0,1)$, the algorithm decodes the sketch to compute a $(O(\phi \log n), \phi)$-ED of $G$ with high probability, in $\otil(n)$ space and $2^{O(n)}$ time.
\end{mainthsimple}

\noindent
We note that at least $\Omega(n\log n)$ space is needed for any small enough $\phi$: for example, if the input graph is a matching of size $n/10$, say, its ED gives a $1-O(\phi \log n)$ fraction of the matching edges. 

The decoding time of our sketch can be made polynomial, at the expense of some loss in the quality of the expander decomposition (similarly to~\cite{streamexpdec}), but keeping the space independent of the sparsity $\phi$.

\begin{mainthsimple}[ED algorithm -- polynomial time decoding]
	\label{cor:mainpoly}
	Let $G=(V,E)$ be a graph given in a dynamic stream. Then, there is an algorithm that maintains a linear sketch of $G$ in $n^{1+o(1)}$ space. For any $\phi \in (0,1)$, the algorithm decodes the sketch to compute a $(\phi \cdot n^{o(1)}, \phi)$-ED of $G$ with high probability, in $n^{1+o(1)}$ space and $\poly(n)$ time.
\end{mainthsimple}

\noindent
In this case we are off by subpolynomial factors in both quality and space complexity as compared to the optimal ones. The actual theorem that we prove allows one to trade-off the loss in quality and increase in space.

\subsubsection*{Lower bound}
Most algorithmic applications of EDs, including the ones mentioned above, do not use just one ED of the input graph. Rather, they use an ED sequence obtained by repeatedly computing an ED of the inter-cluster edges from the previous level. This can be done in two natural ways: by contracting the clusters of an ED (we call this variant CED, for “contraction”), or by removing the intra-cluster edges without changing the vertex set (we call this variant RED, for “removal”). These approaches lead to different results and are used for different applications (e.g.,~\cite{boundlink,mincutjli} contract the clusters, and~\cite{stspectral, maxflow} recurse on inter-cluster edges). In the sequential setting, both CEDs and REDs can be obtained straightforwardly given an ED algorithm. However, this is not so obvious in the streaming model.

One the one hand, one should be able get a sparsity-independent algorithm for computing CEDs in dynamic streams via \cref{cor:main} or \cref{cor:mainpoly}: observe that contracting the vertices of a sparsifier of the input graph $G$  gives a sparsifier of the graph obtained by contracting vertices in~$G$, so the idea would be to maintain an independent copy of our algorithm for each of the $O(\log n)$ levels and contracting vertices in the sparsifier based on the decomposition of the previous level. On the other hand, we show a space lower bound for computing REDs, even in insertion only streams, showing that a dependence on $1/\phi$ is necessary.

\begin{mainthsimple}[RED lower bound]
	\label{th:lbsimple}
	Let $\epsilon,\phi \in (0,1)$ such that $1/n \ll \phi \ll 1/\pylog n$ and $ \epsilon =\otil( \phi)$. Any streaming algorithm that with constant probability computes at least two levels of an $(\epsilon,\phi)$-RED requires \smash{$\widetilde{\Omega}(n/\phi)$} bits of space.
\end{mainthsimple}

\noindent
The result seems to challenge the intuition that EDs become weaker as $\phi$ and $\epsilon$ decrease (note that when $\phi$ is, say, $1/n^2$, an ED can simply consist of connected components). The questions of~(1)~whether this bound can be improved,~(2)~how it scales with the number of levels of RED we compute, and~(3)~whether there are algorithms matching such bounds, remain open.

\subsection{Basic notation}
\label{subsec:basic}

\paragraph{Graph streaming.} In this paper, we will be mostly working in the dynamic graph streaming model, where we know the vertex set $V=[n]$, and we receive a stream of insertions and deletions for undirected edges over $V$. In insertion-only graph streams, the only difference is that previously inserted edges cannot be deleted. At the end of the stream, the graph $G=(V,E)$ consists of the edges that have been inserted and not deleted, and we say that $G$ is given in a (dynamic) stream. When we consider a graph given in a (dynamic) stream, we are implicitly assuming it to have $n$ vertices, without introducing the parameter $n$ explicitly. Also, when we refer to $G$ and $n$ without reintroducing them we are implicitly considering the graph resulting from the input stream and its number of vertices.

A powerful tool for dynamic graph streams is linear sketching, introduced in the seminal work of Ahn, Guha, and McGregor~\cite{linmeas}. The idea is to left-multiply the $\binom{n}{2} \times n$ incidence matrix of the graph by a random $k \times \binom{n}{2}$ matrix for $k \ll \binom{n}{2}$. Since the sketch consists of linear measurements, it automatically handles the case of dynamic streams. We will not be using sketching techniques directly, but rather employ existing algorithms that do. In this paper we restrain ourselves to unweighted edge streams. One may study the same problem in general turnstile graph streams~\cite{weightedlinearspars}, but we do not do this here.

\paragraph{Cuts, volumes, expanders.} Given an unweighted graph $G=(V, E)$, possibly with multiple self-loops on the vertices, we will operate with weighted graphs that approximate $G$ in an appropriate sense. We write $G'=(V',E', w)$ to denote a weighted graph, possibly with multiple weighted self-loops on the vertices. We describe next some notation for such weighted graph $G'$. The same notation carries over to the unweighted graph $G$ by implicitly setting $w$ to assign a weight of $1$ to all edges and self-loops.

For any $A,B \subseteq V'$ we denote by $E'(A,B)$ the edges in $E'$ with one endpoint in $A$ and one in $B$, and by $w(A,B)$ the total weight of edges in $E'(A,B)$. The volume of a cut $S \subseteq V'$ is the sum of the (weighted) degrees, including self-loops, of its vertices. We denote it by $\vol_{G'}(S)$. The sparsity of a cut $\emptyset \neq S \subsetneq V'$ is defined as
\begin{equation*}
	\Phi_{G'}(S) = \frac{w(S,V' \setminus S)}{\min\{\vol_{G'}(S), \vol_{G'}(V' \setminus S)\}} \, .
\end{equation*}
For $\psi \in (0,1)$, we make a distinction between cuts having sparsity less than $\psi$, which we call $\psi$-sparse, and cuts having sparsity at least $\psi$, which we call $\psi$-expanding. The sparsity of $G'$ is the defined as
\begin{equation*}
	\Phi_{G'} = \min_{\emptyset \neq S \subsetneq V'} \Phi_{G'}(S) \, ,
\end{equation*}
and we call $G'$ a $\psi$-expander if all its cuts are $\psi$-expanding, i.e. $\Phi_{G'} \ge \psi$.

For a pair of vertices $e=\{u,v\} \in \binom{V'}{2}$, we denote by $\lambda_{e}(G')$ the edge connectivity of $e$ in $G'$, i.e.
\begin{equation*}
	\lambda_{e}(G') = \min_{\emptyset \neq S \subsetneq V'\, : u \in S, v \in V' \setminus S} w(S,V'\setminus S) \, .
\end{equation*}

\paragraph{Expander decomposition.} As we will treat expander decompositions for the input graph $G$ only, we conveniently use the following additional notation. For a cluster $U \subseteq V$, i.e. a subset of the vertices, and a cut $S \subseteq U$, which is also a subset of the vertices, we denote the number of edges crossing $S$ in $U$ by $\uncut{S}{U}$, i.e. $\uncut{S}{U}=|E(S,U \setminus S)|$. This is the \emph{local} cut of $S$ in $U$. If $U = V$, we use a shorthand notation $\uncutg{S} \coloneqq \uncut{S}{V}$. We call such a cut the \emph{global} cut of $S$, since for $S \subseteq U$ we will be interested in both $\uncut{S}{U}$ and $\uncutg{S}$. Moreover, we drop the subscript from the volume and simply write $\vol(\cdot)$ instead of $\vol_G(\cdot)$. Then, EDs can be defined as follow.

\begin{definition}[Expander decomposition]
	\label{def:expdecclassic}
	Let $G=(V,E)$ and let $\epsilon, \phi \in (0,1)$. A partition $\mathcal{U}$ of~$V$ is an $(\epsilon, \phi)$-expander decomposition (for short, $(\epsilon, \phi)$-ED) of $G$ if
	\begin{enumerate}
		\item \label{property:crossingclassic} $\frac{1}{2}\sum_{U \in \mathcal{U}} \uncutg{U} \le \epsilon |E|$, and
		\item \label{property:expanderclassic} for every $U \in \mathcal{U}$, one has that $G[U]$ is a $\phi$-expander.
	\end{enumerate}
\end{definition}
